% Part: sets-functions-relations
% Chapter: sets
% Section: subsets
% Hindi translation (hi-Deva-IN) of Open Logic commit 9620cc73f9c8e0ad003c514a5d3748f29611c4c0.

\documentclass[../../../include/open-logic-section]{subfiles}

\begin{document}

\olfileid[hi]{sfr}{set}{sub}
\olsection{उपसमुच्चय और घात समुच्चय}

\begin{explain}
हमें प्रायः समुच्चयों की तुलना करनी होगी। ऐसी एक स्वाभाविक तुलना इस प्रकार
की जा सकती है: \emph{एक समुच्चय की प्रत्येक वस्तु दूसरे समुच्चय में भी है}।
यह स्थिति इतनी महत्त्वपूर्ण है कि इसके लिए हम एक नया संकेतन प्रस्तुत करेंगे।
\end{explain}

\begin{defn}[उपसमुच्चय]
यदि समुच्चय~$A$ का प्रत्येक !!{element}, $B$ का भी !!a{element} है, तो हम
कहते हैं कि $A$, $B$ का \emph{उपसमुच्चय} है और $A \subseteq B$ लिखते हैं।
यदि $A$, $B$ का उपसमुच्चय नहीं है, तो हम $A \not\subseteq B$ लिखते हैं।
यदि $A \subseteq B$ किंतु $A \neq B$, तो हम $A \subsetneq B$ लिखते हैं और
कहते हैं कि $A$, $B$ का \emph{उचित~उपसमुच्चय} है।
\end{defn}

\begin{ex}
प्रत्येक समुच्चय अपना ही उपसमुच्चय है और $\emptyset$ हर समुच्चय का
उपसमुच्चय है। सम प्राकृत संख्याओं का समुच्चय, प्राकृत संख्याओं के समुच्चय का
उपसमुच्चय है। साथ ही, $\{ a, b \} \subseteq \{ a, b, c \}$। किंतु
$\{ a, b, e \}$, $\{ a, b, c \}$ का उपसमुच्चय नहीं है।
\end{ex}

\begin{ex}
संख्या~$2$ पूर्णांकों के समुच्चय का !!a{element} है, जबकि सम संख्याओं का
समुच्चय पूर्णांकों के समुच्चय का उपसमुच्चय है। फिर भी कोई समुच्चय किसी दूसरे
समुच्चय का !!a{element} और उपसमुच्चय \emph{दोनों} हो सकता है; उदाहरणतः,
$\{0\} \in \{0, \{0\}\}$ और $\{0\} \subseteq \{0, \{0\}\}$ भी।
\end{ex}

विस्तारिता समुच्चयों की समानता का एक निकष देती है: $A = B$ तभी और केवल तभी
जब $A$ का प्रत्येक !!{element}, $B$ का भी !!a{element} हो, और इसके विपरीत।
``उपसमुच्चय'' की परिभाषा $A \subseteq B$ को इस निकष के ठीक पहले अर्धांश के
रूप में परिभाषित करती है: $A$ का प्रत्येक !!{element}, $B$ का भी
!!a{element} है। निस्संदेह, यही परिभाषा $A$ और $B$ को परस्पर बदलने पर भी
लागू होती है: अर्थात् $B \subseteq A$ तभी और केवल तभी जब $B$ का प्रत्येक
!!{element}, $A$ का भी !!a{element} हो। और यह विस्तारिता का ठीक वही
``इसके विपरीत'' वाला अंश है। दूसरे शब्दों में, विस्तारिता से निष्पन्न होता है
कि दो समुच्चय तभी और केवल तभी समान हैं जब वे एक-दूसरे के उपसमुच्चय हों।

\begin{prop}
$A = B$ तभी और केवल तभी जब $A \subseteq B$ और $B \subseteq A$ दोनों हों।
\end{prop}

अब कुछ और उपयोगी संकेतनों को प्रस्तुत करने का भी अच्छा अवसर है। जब हमने यह
परिभाषित किया कि $A$, $B$ का उपसमुच्चय कब है, तब कहा कि ``$A$ का प्रत्येक
!!{element} \dots'' और ``$\dots$'' के स्थान पर ``$B$ का भी !!a{element}
है'' रखा। किंतु अभिव्यक्ति का यह \emph{रूप} इतना सामान्य है कि इसके लिए एक
औपचारिक संकेतन प्रस्तुत करना उपयोगी होगा।

\begin{defn}\ollabel{forallxina}
$(\forall x \in A)\phi$, $\forall x(x \in A \lif \phi)$ का संक्षिप्त रूप
है। इसी प्रकार, $(\exists x \in A)\phi$, $\exists x(x \in A \land \phi)$
का संक्षिप्त रूप है।
\end{defn}

इस संकेतन का उपयोग करके हम कह सकते हैं कि $A \subseteq B$ तभी और केवल तभी
जब $(\forall x \in A)x \in B$।

अब हम एक विशेष प्रकार के समुच्चय पर विचार करेंगे: किसी दिए हुए समुच्चय के सभी
उपसमुच्चयों का समुच्चय।

\begin{defn}[घात समुच्चय]
समुच्चय~$A$ के सभी उपसमुच्चयों से बने समुच्चय को $A$ का \emph{घात समुच्चय}
कहते हैं और इसे $\Pow{A}$ लिखते हैं।
  \[
    \Pow{A} = \Setabs{B}{B \subseteq A}
  \]
\end{defn}

\begin{ex}
$\{ a, b, c \}$ के सभी संभव उपसमुच्चय कौन-से हैं? वे हैं:
$\emptyset$, $\{a \}$, $\{b\}$, $\{c\}$, $\{a, b\}$, $\{a, c\}$,
$\{b, c\}$, $\{a, b, c\}$। इन सभी उपसमुच्चयों का समुच्चय
$\Pow{\{a,b,c\}}$ है:
\[
\Pow{\{ a, b, c \}} = \{\emptyset, \{a \}, \{b\}, \{c\}, \{a, b\},
\{b, c\}, \{a, c\}, \{a, b, c\}\}
\]
\end{ex}

\begin{prob}
$\{a, b, c, d\}$ के सभी उपसमुच्चयों की सूची बनाइए।
\end{prob}

\begin{prob}
दिखाइए कि यदि $A$ के $n$ !!{element} हैं, तो $\Pow{A}$ के $2^n$
!!{element} हैं।
\end{prob}

\end{document}
